
\documentclass{article}
\usepackage{colortbl}
\usepackage{makecell}
\usepackage{multirow}
\usepackage{supertabular}

\begin{document}

\newcounter{utterance}

\centering \large Interaction Transcript for game `cladder', experiment `full\_v1.5\_default', episode 9560 with selfrefine{-}llama{-}3.1{-}8b{-}instruct.
\vspace{24pt}

{ \footnotesize  \setcounter{utterance}{1}
\setlength{\tabcolsep}{0pt}
\begin{supertabular}{c@{$\;$}|p{.15\linewidth}@{}p{.15\linewidth}p{.15\linewidth}p{.15\linewidth}p{.15\linewidth}p{.15\linewidth}}
    \# & \multicolumn{2}{c}{Player} && \multicolumn{2}{c}{Game Master} \\
    \hline

    \theutterance \stepcounter{utterance}  
    & & & \multicolumn{4}{p{0.6\linewidth}}{
        \cellcolor[rgb]{0.9,0.9,0.9}{
            \makecell[{{p{\linewidth}}}]{
                \texttt{\tiny{[P1$\langle$GM]}}
                \texttt{You are an expert in causal inference. The following question is not a typical commonsense query, but rather a meticulously designed question created by a professor specializing in causal inference, intended to assess the students' mastery of the course content.} \\
\texttt{Imagine a self{-}contained, hypothetical world with only the following conditions, and without any unmentioned factors or causal relationships: Tijv has a direct effect on xevo and gyzp. Xevo has a direct effect on gyzp. We know that tijv causes not xevo. tijv and xevo causes gyzp. We observed an individual is not tijv. Would an individual is gyzp if not xevo instead of xevo?} \\
\texttt{Start your answer with 'Yes' or 'No', followed by additional reasoning or evidence to support your explanation.} \\
            }
        }
    }
     \\ \\

    \theutterance \stepcounter{utterance}  
    & \multicolumn{4}{p{0.6\linewidth}}{
        \cellcolor[rgb]{0.561,0.737,0.561}{
            \makecell[{{p{\linewidth}}}]{
                \texttt{\tiny{[P1$\rangle$GM]}}
                \texttt{FORMAT: ok} \\
\texttt{RULES: ok} \\
\texttt{STRATEGY: problem} \\
\texttt{REASON: The answer assumes the absence of Xevo would not lead to Gyzp, but this conclusion is not directly supported by the given causal relationships.} \\
\\ 
\texttt{Yes} \\
\\ 
\texttt{To answer this question, we need to carefully analyze the given causal relationships and apply the principles of causal inference.} \\
\\ 
\texttt{First, we know that Tijv has a direct effect on Xevo and Gyzp, and Xevo has a direct effect on Gyzp. This establishes a causal pathway from Tijv to Gyzp through Xevo.} \\
\\ 
\texttt{However, we also know that Tijv causes Xevo, which means that the effect of Tijv on Gyzp is mediated by Xevo. In other words, the causal effect of Tijv on Gyzp is indirect and goes through Xevo.} \\
\\ 
\texttt{Now, we are given that an individual is not Tijv. This means that the direct effect of Tijv on Gyzp is eliminated, but the effect of Xevo on Gyzp remains.} \\
\\ 
\texttt{The question asks whether an individual would be Gyzp if not Xevo instead of Xevo. To answer this, we need to consider the counterfactual scenario where Xevo is not present.} \\
\\ 
\texttt{In this scenario, the only remaining causal pathway to Gyzp is the direct effect of Tijv on Gyzp, which is eliminated since the individual is not Tijv. However, we also know that Xevo causes Gyzp, and the individual is not Xevo. } \\
\\ 
\texttt{Since the individual is not Xevo, the direct effect of Xevo on Gyzp is eliminated. However, we also know that Tijv causes Xevo, and the individual is not Tijv. This means that the effect of Tijv on Xevo is eliminated, but it does not directly affect the effect of Xevo on Gyzp.} \\
\\ 
\texttt{Therefore, the answer is No.} \\
            }
        }
    }
    & & \\ \\

    \theutterance \stepcounter{utterance}  
    & & & \multicolumn{2}{p{0.3\linewidth}}{
        \cellcolor[rgb]{0.9,0.9,0.9}{
            \makecell[{{p{\linewidth}}}]{
                \texttt{\tiny{[GM$|$GM]}}
                \texttt{ParseError: The response doesn't start with one of the required word {[}'yes', 'no'{]}, but FORMAT} \\
            }
        }
    }
    & & \\ \\

    \theutterance \stepcounter{utterance}  
    & & & \multicolumn{2}{p{0.3\linewidth}}{
        \cellcolor[rgb]{0.9,0.9,0.9}{
            \makecell[{{p{\linewidth}}}]{
                \texttt{\tiny{[GM$|$GM]}}
                \texttt{game\_result = ABORT} \\
            }
        }
    }
    & & \\ \\

\end{supertabular}
}

\end{document}
